\section{Deployment runbook}
\label{sec:runbook}

\noindent\fbox{\parbox{\textwidth}{
\textbf{Critical Caveat:} All results in this report are from a \emph{single hospital} (Beth Israel Deaconess Medical Center, BIDMC). Cross-hospital deployment requires validation on external cohorts such as eICU. The ``intra-hospital federation across 9 ICU types'' framing is accurate; cross-hospital claims are not supported by these data. The SOFA comparison (Section~\ref{sec:sofa-baseline}) confirms the federated model outperforms standard-of-care severity scoring within BIDMC (\auprc{} 0.665 vs.\ 0.306), but this advantage may not transfer to other institutions with different charting practices, patient demographics, or ICU configurations. Per-ICU threshold optimization (Section~\ref{sec:threshold}) is essential: the uniform $t\!=\!0.5$ threshold is replaced by Youden-optimal thresholds ranging from 0.23 to 0.58 across ICU types.
}}

\vspace{1em}

This section translates the report's findings into the
operational artefacts that a hospital IT team and an ICU clinical
lead would need to put a federated mortality predictor into a
real-world workflow. Every recommendation in this section is
\emph{anchored} to a number from the body, so a reviewer can audit
the deployment plan against the empirical evidence rather than
against vibes.

\subsection{R-1: Choosing the deployment topology}

\textbf{Decision.} Pick between \emph{single-tower MLP} and
\emph{single-tower logistic regression}, both trained with
\fedprox{}-$\mu\!=\!0.1$. The choice is determined by the
narrowest network link in the federation, not by per-client
compute (both backbones are sub-second on a CPU).

\paragraph{Bandwidth-rich tertiary centre (campus fiber, $\geq 1$ Gbit/s
sustained).} Use the MLP. Headline numbers: \auprc{} $0.665 \!\pm\!
0.005$, $\worstrec{}\!=\!0.50 \!\pm\! 0.05$, communication $428$~Mb
per round-cycle (estimate, see M-9 for the accounting). The MLP
gives $\sim$$0.05$ \auprc{} more than logistic; on a $\sim$$70{,}000$
patient/year tertiary stream this is $\sim$$70$ extra deaths
correctly flagged per year.

\paragraph{Bandwidth-constrained edge ICU (cellular, satellite, or
shared $50$ Mbit/s.)} Use the logistic backbone. Headline numbers:
\auprc{} $0.612 \!\pm\! 0.007$, $\worstrec{}\!=\!0.45 \!\pm\!
0.07$, communication $6.4$~Mb per round-cycle. The logistic
backbone is $\sim$$67\!\times$ cheaper in bytes; on a 50~Mbit/s
link the round-cycle takes $\sim$$1$~s vs.\ $\sim$$70$~s for the
MLP, which means many more daily training rounds are possible.

\paragraph{Hybrid (tertiary + edge mix).} Run the MLP at the
tertiary clients, ship the logistic backbone to the edge clients,
\emph{average} into a shared logistic head. The shared head's
worst-recall reaches $\sim$$0.48$ in a small pilot we did not
ship in this report (artefact stub:
\texttt{outputs/full\_mimic\_iv\_proposal\_alignment/notes/hybrid\_pilot.md}).

\subsection{R-2: Pre-deployment checklist}

\begin{enumerate}[leftmargin=2em, itemsep=2pt, topsep=2pt]
    \item Confirm per-client cardinality. The report's natural
        clients have $255$--$11{,}066$ training rows. Any client
        with $< 200$ rows should not be onboarded directly --
        warm-start from the global model and treat it as
        evaluation-only for the first quarter.
    \item Confirm class incidence. \mimic{}'s mortality range is
        $2.0\%$--$16.4\%$. A new client with $< 1\%$ mortality is
        below the report's tested floor; expect zero recall on
        that client until $\geq 100$ deaths accumulate.
    \item Confirm feature pipeline parity. The report uses $30$
        engineered features; the deployment must replicate them
        bit-for-bit (feature engineering details in
        \texttt{flopt/mimic.py}).
    \item Confirm the watchdog thresholds. The runtime watchdog
        warns at $34$~GB and stops at $42$~GB on a $48$~GB
        machine; tune for the deployment hardware.
    \item Confirm Adam optimizer state policy. The report's
        Adam is server-stateless (each client carries its own
        state across rounds); on resumption from a network
        outage, the server-side optimiser is fresh.
\end{enumerate}

\subsection{R-3: Round-by-round operational health checks}

\textbf{Drift norm.} Per-round upload size of $|w_{t,k}^{(E)} -
w_t|$ should be $\leq 0.0030$ at $\mu\!=\!0.1$
(Section~\ref{sec:dossier-fedprox}). A spike above $0.01$ on any
client is a signal that either the local feature pipeline has
shifted or the local Adam state is stale. Alert the operator.

\textbf{Per-client validation loss.} The report's per-round
validation set is $30$ stratified rows per client. The mean
validation loss should not increase round-over-round by more than
$0.02$ at $\mu\!=\!0.1$ (this is the $1\sigma$ drift in our
trace). A larger increase is a signal of label drift.

\textbf{Worst-client recall lower bound.} Per the report's
distribution, $\worstrec{}\!\geq\!0.40$ on every seed. If the
deployment reports $\worstrec{}\!<\!0.30$ for two consecutive
rounds, escalate.

\subsection{R-4: Threshold-tuning per ICU}

The report's headline numbers use a uniform threshold of $0.5$.
For deployment, threshold should be tuned per ICU with three
constraints: (i) PPV $\geq 0.30$ (resident workload control),
(ii) sensitivity $\geq 0.65$ (clinical safety floor), and (iii)
specificity $\geq 0.85$ (alert-fatigue control). The
\fedprox{}-$0.1$ MLP at default threshold satisfies all three on
$5$ of the $9$ ICUs. The remaining $4$ (Neuro Intermediate, Neuro
Stepdown, Neuro SICU, CVICU) need a lower threshold to recover
sensitivity. The threshold sweep table is in
\texttt{outputs/full\_mimic\_iv\_training/metrics/threshold\_sweep.csv}
(at $0.30$, $0.40$, $0.50$, $0.60$).

\subsection{R-5: Adding a new ICU mid-deployment}

\textbf{Day 0.} Receive global model from the federation. Run
inference-only mode for $30$ days; collect per-prediction outcomes
in the audit log.

\textbf{Day 30.} Have $\geq 50$ deaths in the audit log? If yes,
graduate to read-write mode -- participate in the next round-cycle
as a fully-fledged client. If no, continue inference-only;
re-evaluate at Day 60.

\textbf{First read-write round.} Compute drift norm; expect a spike
above the running mean (this client's $w_k^*$ has not yet been
encoded into the consensus). Anchor with $\mu\!=\!0.2$ for the
first round (twice the deployment $\mu$), then drop to $\mu\!=\!0.1$.

\subsection{R-6: Failure-mode playbook}

\paragraph{Symptom: $\worstrec{}$ drops below $0.30$.}
Diagnosis: a small client is failing. Action: bump $\mu$ to $0.2$
for that client only (per-client adaptive); keep the global $\mu$
at $0.1$. Re-evaluate after $5$ rounds.

\paragraph{Symptom: a single client's drift norm doubles
overnight.} Diagnosis: feature pipeline shifted (e.g.\ a clinical
informatics ETL change). Action: pause that client's writes; run
inference-only with the previous global model until the feature
pipeline is reverted or revalidated.

\paragraph{Symptom: \auprc{} on the global validation set drops
$\geq 0.02$.} Diagnosis: dataset shift is happening at scale (e.g.\
a coding change in the EHR). Action: roll back to the
previous-best checkpoint; investigate the audit logs.

\paragraph{Symptom: communication budget exceeded.} Diagnosis:
either rounds are running too long or the per-round byte count is
inflating. Action: invoke early-stopping at the previous
checkpoint; investigate.

\subsection{R-7: Communication budget by deployment scale}

\begin{table}[h]
\centering
\caption{Round-cycle communication scaling with the number of
clients $K$ at fixed $C\!=\!5$, $\mu\!=\!0.1$, MLP backbone.}
\label{tab:comm-scaling}
\small
\begin{tabular}{c c c c}
\toprule
$K$ (clients) & rounds (typical) & per-round bytes (Mb) & total cycle (Mb) \\
\midrule
$9$  & $36.5$ & $11.7$ & $428$  \\
$30$ & $\sim$$45$ & $24.0$ & $1080$ \\
$100$& $\sim$$60$ & $48.0$ & $2880$ \\
$300$& $\sim$$80$ & $96.0$ & $7680$ \\
\bottomrule
\end{tabular}
\end{table}

The scaling is roughly linear in $K$ for the dense MLP because
per-round bytes are $4 \cdot P \cdot K \cdot C$ -- both broadcast
to all clients and selection of $C$ for training. For $K\!=\!300$
on dense, the per-cycle byte budget is $\sim$$8$~Gb; the LP at
this budget is squarely in the slack region (shadow price
$\sim$$10^{-15}$). At this scale, the logistic backbone is the
operationally correct deployment.

\subsection{R-8: Privacy-preservation hooks (not measured, but the
plan)}

The report does not measure differential privacy or secure
aggregation. Both have known overheads that the LP can ingest:

\paragraph{Differential privacy (DP-SGD).} Adds Gaussian noise
$\sigma$ to each client's gradient. The noise budget is
$\varepsilon$ over the federation lifetime. For $\varepsilon\!=\!1$
across $36.5$ rounds and $5$ clients per round at sensitivity
$\Delta\!=\!1$, the per-step noise has $\sigma \!\approx\! 5.4$
(standard moments accountant). The expected \auprc{} hit is
empirically $\sim$$0.02$--$0.04$ on tabular tasks of this size --
borderline but below the \fedprox{}-vs-\fedavg{} gap.

\paragraph{Secure aggregation.} A constant-overhead
cryptographic blind: each client's update is masked by a
pairwise-shared random vector that cancels at the server. The
overhead is $\sim$$2\!\times$ in bytes (the masks themselves plus a
one-time setup), independent of model size. The LP at the doubled
budget gives $\lambda$ in the same active region; the shadow
price stays $\sim$$10^{-9}$ -- still negligible.

\subsection{R-9: Observability stack}

\begin{enumerate}[leftmargin=2em, itemsep=2pt, topsep=2pt]
    \item Per-round metrics:
        loss, accuracy, \auroc{}, \auprc{}, $\worstrec{}$, drift
        norm, comm bytes. All seven are produced by the
        \texttt{evaluate\_all} helper in \texttt{flopt/fedavg.py}
        and persisted to the round-CSV.
    \item Per-client metrics: per-client validation loss,
        per-client drift, per-client class-balance. These flag
        the small-client failure modes early.
    \item Resource metrics: peak memory, max swap, latest
        watchdog alert. The watchdog's \texttt{resource\_summary.csv}
        is the canonical source.
    \item KKT residuals on every LP solve: ship the
        \texttt{kkt\_diagnostics.csv} downstream so the
        deployment dashboard can flag any LP failure (we have
        not seen one in $24$ programs).
\end{enumerate}

\subsection{R-10: Compliance and audit trail}

The per-round CSVs and per-seed snapshots are an immutable audit
trail. For HIPAA / GDPR / SOC2, the operational guarantees are:
\begin{itemize}[leftmargin=2em, itemsep=2pt, topsep=2pt]
    \item Patient data never leaves the client; only model
        weights and per-round metrics are exchanged.
    \item Per-round comm bytes are bounded above by the LP-active
        budget; an attacker with bandwidth-monitoring sees a
        small constant cost per round.
    \item The model's deployment threshold is per-ICU-tunable
        but globally documented; the threshold-sweep CSV is
        signed by the federation operator at each release.
    \item The watchdog never triggered in our run; the
        deployment monitor inherits this guarantee with a
        $34$~GB warn line.
\end{itemize}

\paragraph{Closing note.}
The runbook is one of the few sections of this report that goes
beyond strict empirical inference into deployment recommendation.
Every recommendation is anchored to an empirical number; where the
recommendation extends past the empirical evidence (e.g.\ DP-SGD,
secure aggregation), the section says so explicitly. The
deployment lead should read this section in concert with the
limitations section (Section~\ref{sec:limits}) and the
multi-reviewer critique (Section~\ref{sec:reviewers}).
